\usepackage{geometry}
\geometry{
    a4paper,
    total={170mm,257mm},
    left=15mm,
    right=15mm,
    top=20mm,
    bottom=20mm
}
% --- PACCHETTI BASE ---
\usepackage[utf8]{inputenc}
\usepackage[T1]{fontenc}
\usepackage[italian]{babel}
\usepackage{amsmath, amsfonts, amssymb, amsthm} % Il "must-have" per la matematica
\usepackage{mathtools} % Estensioni per amsmath
\usepackage{xcolor} % Per l'uso dei colori
\usepackage{hyperref} % Per i link cliccabili nell'indice e nel testo
\usepackage{enumitem} % Per liste personalizzate

% --- DEFINIZIONE TEOREMI ---
% Stile con testo in corsivo
\newtheorem{theorem}{Teorema}[section]
\newtheorem{proposition}[theorem]{Proposizione}
\newtheorem{lemma}[theorem]{Lemma}
\newtheorem{corollary}[theorem]{Corollario}

% Stile con testo normale
\theoremstyle{definition}
\newtheorem{definition}[theorem]{Definizione}
\newtheorem{example}[theorem]{Esempio}
\newtheorem{remark}[theorem]{Osservazione}

% --- SCORCIATOIE (COMMANDS) ---
% Insiemi numerici classici
\newcommand{\R}{\mathbb{R}}
\newcommand{\N}{\mathbb{N}}
\newcommand{\Z}{\mathbb{Z}}
\newcommand{\Q}{\mathbb{Q}}
\newcommand{\C}{\mathbb{C}}

% Operatori comuni
\DeclareMathOperator{\dx}{d\!} % Per il differenziale negli integrali
\newcommand{\abs}[1]{\left| #1 \right|} % Valore assoluto automatico
\newcommand{\norm}[1]{\left\| #1 \right\|} % Norma automatica

% --- ESTETICA ---
\hypersetup{
    colorlinks=true,
    linkcolor=blue,
    filecolor=magenta,      
    urlcolor=cyan,
}
\usepackage[most]{tcolorbox}

% Definiamo uno stile per le definizioni
\newtcolorbox{definizione}[1]{
    colback=blue!5!white,    % Sfondo azzurro chiarissimo
    colframe=blue!75!black,  % Bordo blu scuro
    fonttitle=\bfseries,     % Titolo in grassetto
    title=#1,                % Il titolo è il primo argomento
    arc=4mm,                 % Angoli arrotondati
    enhanced,
    attach boxed title to top left={xshift=3mm, yshift=-2mm},
    boxed title style={colback=blue!75!black}
    }
\usepackage{pgfplots}
\pgfplotsset{compat=1.18} % Assicura la compatibilità con le versioni recenti
\usepackage{tikz}
\usepackage{float}
\usepackage{ulem}